\pdfvariable suppressoptionalinfo 611
\pdfvariable trailerid {[<C2622026083100000000000000000000><C2622026083100000000000000000000>]}
\providecommand{\CRevisionRound}{2}
\documentclass[10pt]{article}
\usepackage[a4paper,margin=17mm]{geometry}
\usepackage{amsmath,amssymb,amsthm,booktabs,microtype}
\usepackage[T1]{fontenc}
\usepackage{lmodern}
\usepackage[hidelinks,hypertexnames=false]{hyperref}
\emergencystretch=2em
\newtheorem{theorem}{Theorem}
\newcommand{\R}{\mathbb R}
\newcommand{\tr}{\operatorname{tr}}
\newcommand{\rank}{\operatorname{rank}}
\title{A Complete Floquet--Jordan Atlas\\for the Two-Step Hill Oscillator}
\author{HCS Research Program}
\date{31 August 2026\quad (revision \CRevisionRound)}
\begin{document}
\maketitle
\begin{abstract}
We classify the periodic two-step Hill oscillator for arbitrary real segment
coefficients, including oscillatory, zero-frequency, and hyperbolic pieces.
Entire segment functions give one exact $SL(2,\R)$ monodromy and a closed
Floquet discriminant whose elliptic and hyperbolic regions are complete.
\ifnum\CRevisionRound>0
At the parabolic boundary we distinguish scalar $\pm I$ from nontrivial
Jordan growth, and close every integer iterate and Floquet rate.  An
independent entire-series checker verifies 900 all-sign transfers.
\fi
\ifnum\CRevisionRound>1
All zero-duration, constant-coefficient, and segment-order faces are included.
No arithmetic or target-determinant claim is made.
\fi
\end{abstract}

Hill's periodic-coefficient equation~\cite{Hill1886} and the two-step
piecewise-constant resonance setting~\cite{Golubev1997} provide historical
context.  The all-sign transfer and complete boundary theorem below are
derived directly.

\section{One entire transfer for all segment signs}
Fix the periodic equation
\begin{equation}
 y''+k(t)y=0,\qquad
 k(t)=\begin{cases}k_1,&0\le t<\tau_1,\\
 k_2,&\tau_1\le t<T,\end{cases}
 \quad T=\tau_1+\tau_2>0,                 \label{eq:hill}
\end{equation}
with $k_j\in\R$ and $\tau_j\ge0$.  Define entire functions
\begin{equation}
 C(k,t)=\sum_{m\ge0}\frac{(-k)^mt^{2m}}{(2m)!},\qquad
 S(k,t)=\sum_{m\ge0}\frac{(-k)^mt^{2m+1}}{(2m+1)!}.      \label{eq:CS}
\end{equation}
They obey $C_t=-kS$, $S_t=C$, and $C^2+kS^2=1$.  Hence the state
$Y=(y,y')^{\mathsf T}$ advances through one constant segment by
\begin{equation}
 \Phi(k,t)=\begin{pmatrix}C&S\\-kS&C\end{pmatrix},
 \qquad \det\Phi=1.                                      \label{eq:phi}
\end{equation}
For $k>0$, $k=0$, and $k<0$, this is respectively the cosine--sine,
shear, and cosh--sinh transfer, without a square-root branch convention.

\section{Complete stability and boundary theorem}
Put $M=\Phi(k_2,\tau_2)\Phi(k_1,\tau_1)$ and write $C_j,S_j$ for the
two segment values.  Direct multiplication gives
\begin{equation}
 \det M=1,\qquad
 \Delta=\tr M=2C_1C_2-(k_1+k_2)S_1S_2.                  \label{eq:delta}
\end{equation}

\begin{theorem}\label{thm:atlas}
For the oscillator \eqref{eq:hill}:
\begin{enumerate}
\item If $|\Delta|<2$, the monodromy is elliptic and every solution is
bounded.  Its Floquet angle per unit time is
$T^{-1}\arccos(\Delta/2)$.
\item If $|\Delta|>2$, the multipliers are a real reciprocal pair.  Generic
solutions grow exponentially, one forward eigenline decays, and the positive
growth rate is $T^{-1}\operatorname{arcosh}(|\Delta|/2)$.
\item If $\Delta=\pm2$ and $M=\pm I$, every solution is periodic or
antiperiodic over one coefficient period.  If instead $M\ne\pm I$, then
$\rank(M\mp I)=1$: one Floquet line is bounded and generic solutions grow
linearly, with alternating sign on the minus face.
\end{enumerate}
For every $n\ge1$,
\begin{equation}
 M^n=U_{n-1}(\Delta/2)M-U_{n-2}(\Delta/2)I.              \label{eq:cheb}
\end{equation}
Here $U_{-1}=0$ and $U_0=1$, so \eqref{eq:cheb} includes $n=1$.
\end{theorem}

\paragraph{Proof.}
The characteristic polynomial is $\lambda^2-\Delta\lambda+1$.  Its roots
are distinct unit conjugates for $|\Delta|<2$ and distinct real reciprocals
for $|\Delta|>2$, proving the first two cases.  At trace $\pm2$, a real
$2\times2$ determinant-one matrix is either $\pm I$ or has one Jordan block;
its powers distinguish bounded from linear growth.  Cayley--Hamilton and the
recurrence $U_n(x)=2xU_{n-1}(x)-U_{n-2}(x)$ prove \eqref{eq:cheb}.

\ifnum\CRevisionRound=0
The baseline ends with the exact monodromy and stability theorem.  Executable
receipts and the degenerate-face ledger enter only in later revisions.
\fi

\ifnum\CRevisionRound>0
\section{Exact witnesses and independent receipt}
Trace alone does not settle the parabolic dynamics.  Four exact sentinels are
\[
 I,\quad -I,\quad
 \begin{pmatrix}1&1\\0&1\end{pmatrix},\quad
 \begin{pmatrix}-1&-1\\0&-1\end{pmatrix};              \label{eq:sentinels}
\]
the last two have generic linear growth.  A hyperbolic exact control is
\[
 \Phi(-1,\log2)=
 \begin{pmatrix}5/4&3/4\\3/4&5/4\end{pmatrix},
\]
with multipliers $2$ and $1/2$.

The producer evaluates all $6^2\cdot5^2=900$ parameter rows with
$k_j\in\{-4,-1,0,1,4,9\}$ and
$\tau_j\in\{0,1/4,1/2,1,3/2\}$.  The checker does not call its trigonometric
or hyperbolic implementation; it rebuilds $C,S$ from 90-term power series,
then tests determinant, \eqref{eq:delta}, classification, and every power
through $n=12$.
\begin{center}
\small
\begin{tabular}{lrr}
\toprule
gate & count & result\\
\midrule
all-sign transfer rows & 900 & 19,849 independent assertions\\
exact boundary matrices & 6 & scalar/Jordan/hyperbolic PASS\\
SymPy identities & 289 & determinant/trace/Chebyshev PASS\\
hostile mutations & 41 & 41 rejected\\
\bottomrule
\end{tabular}
\end{center}
Finite rows validate implementation and boundary conventions; the continuum
classification follows from Theorem~\ref{thm:atlas}.
\fi

\ifnum\CRevisionRound>1
\section{Faces, collisions, and Route-A stop}
Equations \eqref{eq:CS}--\eqref{eq:delta} remain valid when either $k_j$ is
zero or negative and when either duration vanishes.  Equal coefficients
collapse to the constant-coefficient transfer over total time.  Swapping the
segments can change $M$, but $\tr(\Phi_2\Phi_1)=\tr(\Phi_1\Phi_2)$, so it
cannot change the Floquet class.

C110 (forced H\'enon Floquet data), C178 (harmonic strobe), C218
(Kelvin--Voigt), C249 (Van der Pol), and C252 (relay oscillator) do not own
this two-step linear Hill theorem.  Earlier ``Hill determinants'' are orbit
Hessians, not equation \eqref{eq:hill}.  This is collision bookkeeping, not
a priority claim.

The strict assessment is
\[
\texttt{(A0\_FAIL,A1\_WEAK,A2\_FAIL,A3\_FAIL,A4\_FORMAL\_HINT)},
\]
with \texttt{ROUTE\_A\_REJECTED}, Route B false, and scope
\texttt{NO\_BAD\_EULER\_OR\_ROOT\_NUMBER}.  The source characteristic
polynomial is not a target determinant.  No arithmetic local data, Euler
factor, root number, automorphy, target divisor or functional equation,
target zero match, or Hilbert--P\'olya operator is claimed.
\fi

\section{Conclusion and limitations}
The all-sign entire transfer, one discriminant, the scalar/Jordan split, and
the Chebyshev law together close the two-step family.  This does not classify
arbitrary periodic coefficients, nonlinear perturbations, or every initial
condition's forward behavior in the hyperbolic case.

\begin{thebibliography}{9}
\bibitem{Hill1886} G.~W.~Hill, ``On the part of the motion of the lunar
perigee which is a function of the mean motions of the sun and moon,''
\emph{Acta Mathematica} 8 (1886), 1--36,
\href{https://doi.org/10.1007/BF02417081}{doi:10.1007/BF02417081}.
\bibitem{Golubev1997} Y.~F.~Golubev, ``Resonances of linear differential
equations with piecewise constant coefficients,'' Keldysh Institute Preprint
43 (1997), \href{https://www.mathnet.ru/eng/ipmp1431}{MathNet record}.
\end{thebibliography}
\end{document}
